\section{A finite reduction of the four-regular problem}\label{sec:four-finite}

This section combines the doubled-edge reduction with a published degree-sequence bound.
It proves that any minimum counterexample has between eleven and eighteen vertices.
It does not decide whether such a counterexample exists.
The upper cutoff uses exact rational arithmetic for finitely many numerical inequalities.
It does not use an exhaustive graph search at those orders.

\subsection{The degree-sequence input}

We use Theorem~7 of Dvo\v r\'ak, Morrison, Noel, Norin, and Postle,
\emph{Bounding the number of cycles in a graph in terms of its degree sequence},
\url{https://arxiv.org/abs/1907.12091}.
Their theorem gives, for every simple graph $F$ and distinct vertices $s,t$,
\[
 P_{st}(F)\le a(\deg_F(s))
                \prod_{v\ne s}b(\deg_F(v)).
\]
Only the following values enter this argument:
\[
 \rho=(4/3)^{1/4},\quad b(2)=\rho,\quad
 b(3)=(2\rho)^{1/2},\quad b(4)=(6\rho)^{1/3}=2/\rho,
 \quad b(5)=(24\rho)^{1/4},
\]
\[
 a(2)=2/\rho^2,\qquad a(3)=6/(\rho^4b(5)).
\]
The theorem permits disconnected graphs.
The product still gives an upper bound when the terminals lie in different components.

Let $G$ be a connected loopless four-regular graph on $n$ vertices.
Assume that every multiplicity is at most two and its $k$ doubled pairs form a matching.
Subdivide one edge from each doubled pair, using a distinct new vertex each time.
The resulting graph $S$ is simple and has exactly the same cycles as $G$.
Its degree product is
\begin{equation}\label{eq:core-degree-product}
 B=b(4)^n\rho^k=2^n\rho^{k-n}
 \le2^n(3/4)^{n/8},
\end{equation}
since $k\le n/2$.

\begin{lemma}\label{lem:core-edge-frequency}
Set
\[
 \kappa=\frac{a(3)b(3)}{b(4)^2}
        =\frac{3^{19/16}}{2^{17/8}}<\frac{423}{500}.
\]
Every edge of $G$ satisfies $c_e(G)\le\kappa B$.
If $e$ belongs to a doubled pair, it also satisfies $c_e(G)\le1+B/2$.
Consequently,
\begin{equation}\label{eq:core-length-sum}
 \sum_C|E(C)|
 \le \bigl(2\kappa(n-k)+k\bigr)B+2k.
\end{equation}
\end{lemma}

\begin{proof}
First suppose that the edge survives unsubdivided in $S$.
Its endpoints have degree four.
Delete this edge and apply the path theorem to its endpoints.
Their degrees become three, so
\[
 c_e(G)=c_e(S)=P_{uv}(S-e)
 \le a(3)b(3)b(4)^{n-2}\rho^k=\kappa B.
\]
The two original edges in a doubled pair have equal cycle counts.
Thus the same bound covers the edge that was subdivided.

Now let $e$ belong to a doubled pair with ends $u,v$.
Delete both edges to obtain $N$.
Every cycle through $e$ is either the two-edge cycle or a path in $N$ closed by $e$.
Subdivide one edge from every remaining doubled pair.
The resulting simple graph has terminal degrees two, $n-2$ original degrees four, and $k-1$ additional degrees two.
Hence
\[
 P_{uv}(N)\le a(2)\rho^k b(4)^{n-2}=B/2.
\]
This proves the second bound.
Sum it over the $2k$ doubled edges.
Use the first bound for the other $2n-2k$ edges.
Each cycle is counted once for every edge it uses, proving \eqref{eq:core-length-sum}.
\end{proof}

\subsection{A short-cycle bound that records the doubled matching}

Write $c_\ell(G)$ for the number of cycles of length $\ell$.
There are exactly $k$ two-edge cycles.
For $\ell\ge3$, define
\begin{equation}\label{eq:core-U}
 U_\ell(n,k)=\min\left\{
  \frac{(6n+2k)3^{\ell-3}}\ell,
  \frac{49}{81\ell}\left(n-\frac{2k}{9}\right)
       \left(\frac{18}{7}\right)^\ell
       \left(\frac76\right)^{n-2k}
 \right\}.
\end{equation}

\begin{lemma}\label{lem:core-short-cycles}
For every $\ell\ge3$, we have $c_\ell(G)\le U_\ell(n,k)$.
\end{lemma}

\begin{proof}
First fix a starting vertex incident with no doubled pair.
There are four first-edge choices, at most three intermediate choices, and at most one closing edge.
This gives at most $4\cdot3^{\ell-2}$ directed cycle representations at that vertex.

Now fix a starting vertex incident with a doubled pair.
Classify its cycles by the two edges used there.
If both are single, there are at most $2\cdot3^{\ell-2}$ directed representations.
If the first edge is doubled, there are two first-edge choices and at most two second-edge choices.
The closing edge must be single.
This gives at most $4\cdot3^{\ell-3}$ representations.
Reversal gives the same bound when only the closing edge is doubled.
Thus each such starting vertex contributes at most $14\cdot3^{\ell-3}$ representations.

There are $n-2k$ vertices of the first type and $2k$ of the second type.
Sum these bounds and divide by $2\ell$.
The result is the first bound in \eqref{eq:core-U}.

For the second bound, call an edge doubled if it belongs to a doubled pair.
Give each directed doubled edge weight $w=7/9$ and each other directed edge weight one.
Put $\lambda=18/7$ and $R=7/6$.
There are $h=n-2k$ vertices incident with no doubled pair.
Call these vertices ordinary.

Consider a vertex incident with a doubled pair.
After arriving along a single edge, a vertex-simple walk has at most one single and two doubled choices.
Their total weight is $1+2w=23/9\le\lambda$.
After arriving along a doubled edge, returning along its parallel mate is forbidden.
At most two single choices remain, with total weight $2=\lambda w$.
At an ordinary vertex, at most three choices remain, and $3=\lambda R$.

For completeness, these weighted inequalities can be multiplied along a branching search.
At each partial vertex-simple walk, record its last-edge weight and the unused number of ordinary vertices.
With $j$ steps remaining, multiply them by $\lambda^jR^h$, using the remaining value of $h$.
The total of the children's quantities never exceeds the parent's quantity.
Each ordinary vertex is charged at most once.

At the final step, require closure to the fixed starting vertex.
After a single arrival at a nonordinary vertex, closing edges have total weight at most $2w=14/9$.
Relative to $\lambda$, this is $49/81$.
After a doubled arrival, a closing edge must be single when $\ell\ge3$.
Its relative factor is at most $1/(\lambda w)=1/2$.
At an ordinary vertex the relative factor is at most $1/3$.
Thus closure supplies a factor at most $49/81$ in every case.

The total starting-edge weight is
\[
 W=(4n-4k)+4kw=4n-8k/9.
\]
The weighted count of closed vertex-simple walks is therefore at most
\[
 \frac{49}{81}W\lambda^{\ell-1}R^h.
\]
Every last edge has weight at least $w$.
Divide by $w$ and then by $2\ell$.
Since $\lambda w=2$, the result is the second bound in \eqref{eq:core-U}.
\end{proof}

\subsection{The explicit cutoff}

\begin{theorem}\label{thm:four-core-cutoff}
Let $G$ be a connected loopless four-regular graph whose doubled edges form a matching and whose other edges are single.
If $|V(G)|\ge19$, then $c(G)<2^{|V(G)|}$.
Consequently, a minimum counterexample to $c(G)\le2^n+n$ has $11\le n\le18$.
\end{theorem}

\begin{proof}
First consider $19\le n\le32$.
For $0\le k\le\lfloor n/2\rfloor$, put
\[
 g=\left\lfloor\frac{4n+3k}{7}\right\rfloor,
 \qquad \bar\kappa=423/500.
\]
For every cycle length $\ell$, its contribution to $c(G)$ is at most $\ell/g$,
except that a shorter cycle requires the additional contribution $1-\ell/g$.
Apply Lemmas~\ref{lem:core-edge-frequency} and~\ref{lem:core-short-cycles}.
The two-edge contributions simplify to $k$, giving
\begin{equation}\label{eq:core-finite-bound}
 \frac{c(G)}{2^n}\le A_{n,k}+T_{n,k}(3/4)^{(n-k)/4},
\end{equation}
where
\[
 A_{n,k}=2^{-n}\left(k+
   \sum_{\ell=3}^{g-1}(1-\ell/g)U_\ell(n,k)\right),
 \qquad T_{n,k}=\frac{2\bar\kappa(n-k)+k}{g}.
\]
Both quantities are rational.
For every stated pair $(n,k)$, exact arithmetic gives
\begin{equation}\label{eq:core-rational-certificate}
 A_{n,k}<1,\qquad
 T_{n,k}^4(3/4)^{n-k}<(1-A_{n,k})^4.
\end{equation}
Thus \eqref{eq:core-finite-bound} is less than one.
The verification script records every rational inequality, rather than testing rounded approximations.
For reference, the largest displayed bounds at each order are below the following rounded-up values:
\[
\begin{array}{c|rrrrrrr}
n&19&20&21&22&23&24&25\\\hline
\text{bound}&.984&.956&.922&.902&.846&.827&.800\\
n&26&27&28&29&30&31&32\\\hline
\text{bound}&.788&.778&.769&.696&.689&.683&.678
\end{array}
\]

For $n\ge33$, a simpler bound gives an induction over all orders.
Put $g=\lfloor4n/7\rfloor$ and $q=(3/4)^{1/8}$.
Use $c_2(G)\le n/2$, the elementary bound $c_\ell(G)\le2n3^{\ell-1}/\ell$,
and $\sum_C|E(C)|\le2n\kappa B$.
Then
\begin{equation}\label{eq:core-tail-bound}
 \frac{c(G)}{2^n}
 \le F_n:=\frac{n/2+2nS_g}{2^n}+
              \frac{2n}{g}\frac{17}{20}q^n,
 \qquad S_g=\sum_{\ell=3}^{g-1}\frac{3^{\ell-1}}\ell.
\end{equation}
Exact rational eighth-power tests show $F_n<1$ for $33\le n\le39$.
These are the seven base cases for the induction.

For $g\ge5$, we have $S_g\ge3^{g-1}/(2g)$.
The claim holds at $g=5$, and its induction follows by adding $3^{g-1}/g$.
It implies $S_{g+1}\le3S_g$, and hence $S_{g+4}\le81S_g$.
Under $n\mapsto n+7$, the cutoff $g$ increases by four.
The first term in \eqref{eq:core-tail-bound} decreases, since
\[
 \frac{81(n+7)}{128n}<1\qquad(n\ge33).
\]
The second term decreases because $7g\le4n$ gives $(n+7)/(g+4)\le n/g$, and $q<1$.
This proves $F_n<1$ for all $n\ge33$.

Finally, Section~\ref{sec:four-regular} reduces a minimum counterexample to the stated class.
The Hamilton-pair enumeration now verifies all orders through ten.
The analytic bound excludes every order at least nineteen.
\end{proof}

\paragraph{Certificates and scope.}
The script \texttt{check\_four\_regular\_core\_bound.py} checks all finite inequalities with rational arithmetic.
It also verifies $\kappa^{16}=3^{19}/2^{34}<(423/500)^{16}$.
Its complete output is \texttt{four-regular-core-bound-results.json}.
The independent graph enumeration through order ten is in \texttt{four-regular-phase2-results.json}.
The new order contributes $181{,}440$ Hamilton-pair orderings.
Its maximum is $1034=2^{10}+10$.
The bounds and reductions above are not formalised in Lean.


\subsection{Sharper cutoffs with at most two doubled pairs}

The short-cycle count also gives useful cutoffs for a fixed small value of $k$.
These do not require a Hamilton decomposition.

\begin{corollary}\label{cor:small-matching-cutoffs}
Under the hypotheses of Theorem~\ref{thm:four-core-cutoff}, we have $c(G)<2^n$ in each of these cases:
\[
 (k=0\text{ and }n\ge14),\qquad
 (k=1\text{ and }n\ge15),\qquad
 (k=2\text{ and }n\ge16).
\]
\end{corollary}

\begin{proof}
Theorem~\ref{thm:four-core-cutoff} covers $n\ge19$.
For the remaining orders, use \eqref{eq:core-finite-bound} with the following cutoffs:
\[
\begin{array}{c|ccccc}
 n&14&15&16&17&18\\\hline
 k=0&10&10&11&12&12\\
 k=1&-&10&11&12&12\\
 k=2&-&-&11&12&12
\end{array}
\]
Each case satisfies the exact rational inequalities in \eqref{eq:core-rational-certificate}.
The largest displayed upper ratio is less than $.974$.
The certificate script records all twelve checks.
\end{proof}

If a separate structural reduction leaves only $k\le2$, these cutoffs leave the following orders to check:
\[
\begin{array}{c|ccc}
 k&0&1&2\\\hline
 n&11,12,13&11,12,13,14&11,12,13,14,15.
\end{array}
\]
This conditional statement uses the exhaustive Hamilton-pair verification through order ten.
It does not assert that these remaining graph classes have already been enumerated.

\subsection{The unresolved contraction inequality}

Let $u,v$ be joined by exactly two edges in a four-regular graph $G$ on $n$ vertices.
Delete these edges to obtain $N$.
Let $H$ identify $u,v$ in $N$.
Thus $H$ is the four-regular graph obtained by contracting the doubled pair and deleting its two loops.
Assume that $N$ is connected.
Write $c_{uv}(N)$ for the number of cycles containing both terminals.
Direct classification gives
\[
 c(G)=c(N)+2P_{uv}(N)+1,
 \qquad c(H)=c(N)-c_{uv}(N)+P_{uv}(N).
\]
Consequently,
\[
 c(G)-c(H)-1=P_{uv}(N)+c_{uv}(N).
\]
Since $\beta(N)=n-1$, the proposed contraction inequality is exactly
\[
 P_{uv}(N)+c_{uv}(N)\le2^{\beta(N)}.
\]
The terminals have degree two, and every other vertex has degree four.
When $G$ has a Hamilton decomposition, $N$ is the union of two Hamilton terminal paths.
The inequality passed all doubled-edge contraction tests through order ten.
It remains unproved and is not used in Theorem~\ref{thm:four-core-cutoff}.
